



\chapter{Classifying (Grouping) Hosts}





\section{Hard classes}





\begin{description}
\item[\textbf{Class}] A group of things, animals, or people with similar features or qualities. ---Macmillan Dictionary
\end{description}

\begin{quote}
Classes are the if ( test ) then of CFEngine language. Tests are built-in
or user defined. Hosts that pass the test are members of the class.
---Neil Watson, CFEngine Consultant
\end{quote}





\subsection{Hard and soft classes defined}

There are two types of classes in CFEngine:

\begin{description}
\item[\textbf{Hard classes}] Hard classes are inherent, or built-in. The first thing that \kode{cf-agent}
does when it starts is to classify its environment (e.g.  OS type = linux,
OS version = redhat 6.5, date = Sun Nov  8 17:09:57 PST 2015, CFEngine
version = 3.7, hostname = alpha.example.com, domain = example.com, CPU
architecture = 64 bit, etc.)  This data can be used to control promise
execution (e.g. kick off backups at 2 AM on Sunday on Linux hosts)



\item[\textbf{Soft classes}] Soft classes are user-defined through promises, and provide additional
flexibility in classifying hosts (e.g. by application, or primary vs
DR) and controlling promise execution (e.g. only evaluate promise2 if
promise1 was repaired).
\end{description}





\subsubsection{Note on Syntax: CFEngine Identifiers}

Note on syntax: CFEngine identifiers (class names, variable names, bundle
names, etc.) may only contain alphanumeric and underscore characters (a-zA-Z0-9\_).





\subsection{Hard classes}

Let's see some examples of hard classes.

\begin{codelisting}
\codecaption{300-010-Basic\_Examples\_Classes-0030-Classes\_Reports.cf}
%= lang:cfengine3, options: "linenos": false
\begin{code}
# Operating System

bundle agent main
{
  commands:
    linux::
      "/bin/date";

    windows::
      "C:\Windows\System32\cmd.exe /c date /t";
}
\end{code}
\end{codelisting}
\begin{codelisting}
\codecaption{300-010-Basic\_Examples\_Classes-0040-Using\_classes\_to\_determine\_role.cf}
%= lang:cfengine3, options: "linenos": false
\begin{code}
# IPv4 network blocks

bundle agent main
{
  reports:
    ipv4_205_186_156::
      "I am on our public net. I am a Web server.";

    ipv4_10::
      "I am on our private net. I am a database server.";
}
\end{code}
\end{codelisting}
\begin{codelisting}
\codecaption{300-010-Basic\_Examples\_Classes-0050-Report\_OS\_Type.cf}
%= lang:cfengine3, options: "linenos": false
\begin{code}
# OS flavor (e.g. Windows XP or Red Hat)

bundle agent main
{
  reports:
    WinXP:: "Hello world! I am running on a Windows system.";
    linux:: "Hello world! I am running on a Linux system.";
    redhat:: "Hello world! I am running on a redhat Linux system.";
}
\end{code}
\end{codelisting}
\begin{codelisting}
\codecaption{300-010-Basic\_Examples\_Classes-0060-Note\_on\_what\_happens\_to\_dashes\_in\_hostnames.cf}
%= lang:cfengine3, options: "linenos": false
\begin{code}
# CFEngine automatically canonifies classes (converts any
# character that is not alphanum/underscore to underscore)
#
# To setup for this example, run "hostname my-hostname-has-dashes"

bundle agent main
{
  reports:
    any::
      "hello world";
    my-hostname-has-dashes::
      "One";
    my_hostname_has_dashes::
      "Two";

}

\end{code}
\end{codelisting}





\begin{aside}
\label{aside:exercise_18}
\heading{Examine hard classes}

Run CFEngine in verbose mode:

%= lang:bash
\begin{code}
cf-agent -v -f ./hello_world.cf | less
\end{code}

Examine what CFEngine discovered about your system and what classes it set.

Give an example of a class that CFEngine has set.

\end{aside}





\begin{aside}
\label{aside:exercise_19}
\heading{Using classes}

Print a report if you're running on a CentOS 7 system.

\end{aside}





\section{Class Expressions}





\begin{description}
\item[\textbf{Class expressions}] Class expressions are logical expressions that evaluate to true (this
promise applies here) or false (the promise does not apply, skip it).
\end{description}

Class expressions are composed of classes and logical operators.





\subsection{Logical operators}

\begin{description}
\item[\textbf{Operator}] In programming, a symbol used to perform an arithmetic or logical operation.
\end{description}

--- http://encyclopedia2.thefreedictionary.com/operator

Logical operators (in order of precedence of operation)

\begin{longtable}{|l|l|}
\hline
( ) & Groupers\\
! & NOT\\
\& or . & AND\\
\textbar{} or \textbar{}\textbar{} & OR\\
\hline
\end{longtable}





\subsection{Truth tables}

If necessary, review \href{https://en.wikipedia.org/wiki/Truth_table#Logical_conjunction_.28AND.29}{truth tables} for logical operations AND, OR, and NOT.





\subsection{Examples of class expressions}

\begin{codelisting}
\codecaption{300-015-Basic\_Examples\_Classes\_2-0090-Class\_expression\_operators.cf}
%= lang:cfengine3, options: "linenos": false
\begin{code}
bundle agent main
{
  reports:
    !WinXP:: "This isn't Windows XP";
    windows|linux::  "Am I laughing or crying?";
    windows&linux::  "We should never see this report.";

}
\end{code}
\end{codelisting}
\begin{codelisting}
\codecaption{300-015-Basic\_Examples\_Classes\_2-0100-Report\_day\_of\_the\_week.cf}
%= lang:cfengine3, options: "linenos": false
\begin{code}
# This bundle does not use class expressions.

bundle agent main
{
  reports:

    Monday::    "Hello world! I love Mondays!";
    Tuesday::   "Hello world! I love Tuesdays!";
    Wednesday:: "Hello world! I love Wednesdays!";
    Thursday::  "Hello world! I love Thursdays!";
    Friday::    "Hello world! I love Fridays!";

    Saturday::  "Hello world! I love weekends!";
    Sunday::    "Hello world! I love weekends!";
}
\end{code}
\end{codelisting}
\begin{codelisting}
\codecaption{300-015-Basic\_Examples\_Classes\_2-0110-Condensed\_report\_day\_of\_week.cf}
%= lang:cfengine3, options: "linenos": false
\begin{code}
# This bundle uses class expressions.

bundle agent main
{
  reports:
    Monday|Tuesday|Wednesday|Thursday|Friday::
      "I get to work today";

    Saturday|Sunday::
      "I get to rest today.";
}
\end{code}
\end{codelisting}
\begin{codelisting}
\codecaption{300-015-Basic\_Examples\_Classes\_2-0120-OS\_and\_time\_expression.cf}
%= lang:cfengine3, options: "linenos": false
\begin{code}
# Another example of a class expression

bundle agent main
{
  reports:
      linux&Hr22::
        "Linux system AND we are in the 22nd hour.";
}

\end{code}
\end{codelisting}
\begin{codelisting}
\codecaption{300-015-Basic\_Examples\_Classes\_2-0130-Class\_expression\_OS\_and\_time.cf}
%= lang:cfengine3, options: "linenos": false
\begin{code}
# Example of using ( ) for grouping

bundle agent main
{
  reports:
    (redhat&Monday)|(windows&Wednesday)::
      "This report will show on Redhat servers on Mondays;
       or on Windows servers on Wednesdays";
}
\end{code}
\end{codelisting}





\section{Soft classes}

You can define a soft class using a \kode{classes:} promise.

\begin{codelisting}
\codecaption{300-015-Basic\_Examples\_Classes\_2-0134-soft\_classes\_simple.cf}
%= lang:cfengine3, options: "linenos": false
\begin{code}
bundle agent main
{
  classes:
      "weekend"
        expression => "Saturday|Sunday";
      "weekday"
        expression => "Monday|Tuesday|Wednesday|Thursday|Friday";

  reports:
    weekend::
      "Yay! I get to rest today.";
    weekday::
      "Yay! I get to work today.";
}
\end{code}
\end{codelisting}
\begin{codelisting}
\codecaption{300-015-Basic\_Examples\_Classes\_2-0140-negative\_knowledge.cf}
%= lang:cfengine3, options: "linenos": false
\begin{code}
# Example of "negative knowledge" -- not recommended!
# Better to be certain (rely on the presence of something,
# not its absence).

bundle agent main
{
  classes:
      "weekend"
        expression => "Saturday|Sunday";
      "weekday"
        not => "weekend";

  reports:
    weekend::
      "Yay! I get to rest today.";
    weekday::
      "Yay! I get to work today.";
}

# Knowing that something is not the case is not the same as not knowing
# whether something is the case. That a class is not set could mean
# either.
#
# Reference:  <https://docs.cfengine.com/docs/3.17/reference-language-concepts-classes.html#negative-knowledge>
\end{code}
\end{codelisting}
\begin{codelisting}
\codecaption{300-030-Basic\_Examples\_Classes\_2-0140-Detect\_VMs.cf}
%= lang:cfengine3, options: "linenos": false
\begin{code}
# You can set a soft class based on the outcome
# of a function that returns true/false, such as
# `regline()` which checks if there is a line in a
# file matching a regular expression.

bundle agent main
{
  classes:
      "i_am_virtual"
        comment => "Check if we are running inside a VM",
        expression => regline(".*(VMware|VBOX|QEMU).*",
                              "/proc/scsi/scsi");

# E.g., on a VMware guest, we have:
#
# $ grep -i vmware /proc/scsi/scsi
# Vendor: VMware,  Model: VMware Virtual S Rev: 1.0
# Vendor: NECVMWar Model: VMware SATA CD01 Rev: 1.00
# $

  reports:
    i_am_virtual::
      "Running inside a VM";
}

# See also the "virt-what" utility
\end{code}
\end{codelisting}





\subsection{Creating soft classes depending on promise outcome}

Some promise attributes can create Classes depending on the outcome of
the promise.

\begin{codelisting}
\codecaption{300-040-Classes\_4-0020-Ensuring\_CUPSd\_is\_running.cf}
%= lang:cfengine3, options: "linenos": false
\begin{code}
# restart_class will set the class if the process is ABSENT
# <https://docs.cfengine.com/latest/reference-promise-types-processes.html#restart_class>

bundle agent main
{
  processes:
    print_servers::
      "cupsd"
        restart_class => "cups_needs_to_be_started",
        comment => "We want print services";

  commands:
    cups_needs_to_be_started::
      "/sbin/service cups start";
}
\end{code}
\end{codelisting}
\begin{codelisting}
\codecaption{300-040-Classes\_4-0030-Ensuring\_httpd\_is\_running.cf}
%= lang:cfengine3, options: "linenos": false
\begin{code}
bundle agent main
{
  processes:
      "httpd"
        restart_class => "start_httpd";

  commands:
    start_httpd::
      "/sbin/service httpd start";
}
\end{code}
\end{codelisting}